\chapter{Introduction}\label{introduction}
\chapterstatus{complete}

\section{What this project is}\label{introduction:section-aim}

The Hodge Project is a cumulative reference whose goal is to develop, with
complete proofs, all the mathematics needed to state, understand and work with
the Hodge conjecture, together with the theory that surrounds it.

Informally, the Hodge conjecture says the following: on a smooth complex
projective variety $X$, every rational cohomology class of type $(p,p)$ is a
rational linear combination of the classes of algebraic subvarieties of $X$. A
precise statement is given in Chapter~\ref{hodge-conjecture}. Each word of the
informal statement hides a substantial theory: rational cohomology classes
come from algebraic topology (Chapter~\ref{singular-cohomology}), the type
$(p,p)$ comes from the Hodge decomposition of a compact K\"ahler manifold
(Chapter~\ref{hodge-decomposition}), which rests on the analysis of elliptic
operators (Chapter~\ref{elliptic-operators}), and the classes of subvarieties
come from algebraic geometry (Chapters~\ref{varieties} and~\ref{cycle-class}).
This project builds all of it, starting from logic and set theory.

The organisation is modelled on the Stacks project: the text is written in a
strict logical order, every result has a permanent tag, and the whole project
is a single document that can be read as a book or browsed as a website.

\section{How the project is organised}\label{introduction:section-organisation}

The chapters are grouped into parts.
\begin{itemize}
\item Part~\ref{part-foundations}, \emph{Foundations}: logic, set theory, the
number systems and category theory.
\item Part~\ref{part-algebra}, \emph{Algebra}: from groups to derived
categories.
\item Part~\ref{part-topology}, \emph{Topology}: general topology, homotopy
theory, singular homology and cohomology, Poincar\'e duality.
\item Part~\ref{part-sheaves}, \emph{Sheaves}: sheaves and sheaf cohomology.
\item Part~\ref{part-analysis}, \emph{Analysis}: real, complex and functional
analysis, several complex variables and Sobolev spaces.
\item Part~\ref{part-differential-geometry}, \emph{Differential Geometry}: from
smooth manifolds to the Hodge theorem for harmonic forms.
\item Part~\ref{part-complex-geometry}, \emph{Complex Geometry}: complex and
K\"ahler manifolds, the Hodge decomposition, the Lefschetz theorems, the
Kodaira vanishing and embedding theorems, and Calabi--Yau manifolds.
\item Part~\ref{part-algebraic-geometry}, \emph{Algebraic Geometry}: varieties,
schemes, coherent cohomology, GAGA, curves, algebraic groups, abelian
varieties, algebraic cycles, the cycle class map and \'etale cohomology.
\item Part~\ref{part-hodge-theory}, \emph{Hodge Theory}: Hodge structures and
Mumford--Tate groups, the Hodge conjecture, the evidence for it, its failed
variants, mixed Hodge structures, absolute Hodge classes, and the Tate and
standard conjectures around it.
\end{itemize}
The order of the chapters is the logical order of the project, not a
recommended reading order. A reader interested mainly in the Hodge conjecture
can start from Part~\ref{part-hodge-theory} and follow references backwards
as needed.

\section{A path to the Hodge conjecture}\label{introduction:section-path}

The following chapters form the spine of the project; everything else supports
them or grows out of them.
\begin{enumerate}
\item \emph{Cohomology.} Singular cohomology and Poincar\'e duality
(Chapters~\ref{singular-cohomology} and~\ref{poincare-duality}), sheaf
cohomology (Chapter~\ref{sheaf-cohomology}) and de Rham's theorem
(Chapter~\ref{de-rham}) give the rational cohomology of a manifold and its
description by differential forms.
\item \emph{Harmonic theory.} Elliptic operators
(Chapter~\ref{elliptic-operators}) give the Hodge theorem: every de Rham class
has a unique harmonic representative (Chapter~\ref{hodge-theorem}).
\item \emph{K\"ahler geometry.} On a compact K\"ahler manifold
(Chapter~\ref{kahler}) the Laplacian respects the bidegree of forms, which gives
the Hodge decomposition $H^k = \bigoplus_{p+q=k}H^{p,q}$
(Chapter~\ref{hodge-decomposition}), hard Lefschetz and the Hodge--Riemann
relations (Chapter~\ref{lefschetz-theorems}).
\item \emph{Algebraic geometry.} Projective varieties are K\"ahler, and GAGA
identifies analytic and algebraic objects on them (Chapter~\ref{gaga}).
Subvarieties have cohomology classes, which are rational and of type $(p,p)$
(Chapter~\ref{cycle-class}).
\item \emph{The conjecture.} Hodge structures (Chapter~\ref{hodge-structures})
isolate the linear algebra, and the Hodge conjecture asserts the converse of
the last statement (Chapter~\ref{hodge-conjecture}). It is a theorem for
divisors (Chapter~\ref{lefschetz-11}), it is known in further cases
(Chapter~\ref{known-cases}), and its natural strengthenings are false
(Chapter~\ref{failures}).
\item \emph{The surrounding theory.} Variations of Hodge structure and normal
functions (Chapters~\ref{vhs} and~\ref{intermediate-jacobians}), mixed Hodge
structures (Chapter~\ref{mixed-hodge-structures}), absolute Hodge classes
(Chapter~\ref{absolute-hodge}), and the Tate and standard conjectures
(Chapters~\ref{tate-conjecture} and~\ref{standard-conjectures}) place the
conjecture in its wider landscape.
\end{enumerate}

\section{Sources and status}\label{introduction:section-sources}

The mathematics in this project is standard, and each chapter lists the
references it follows; the bibliography collects them. The chapters of
Part~\ref{part-hodge-theory} are based in part on the expository text
\emph{The Hodge Conjecture} (2026), whose claims were checked against the
original literature and corrected where necessary. Where a theorem is quoted
rather than proved, the text says so and gives a reference
(Remark~\ref{introduction:remark-proofs}). All chapters are written; none has
yet been checked by an independent reader, which is the condition for the status
\emph{checked} (Remark~\ref{introduction:remark-status}).

\section{Conventions of the text}\label{introduction:section-conventions}

\begin{notation}\label{introduction:notation-numbering}
All numbered items of a section share a single counter: Lemma~5.2.3 is the
third numbered item of Section~5.2, and it is followed by Example~5.2.4.
Chapters are numbered consecutively through all the parts.
\end{notation}

\begin{notation}\label{introduction:notation-environments}
The text uses the following kinds of numbered items.
\begin{itemize}
\item \textbf{Definition}: introduces a notion.
\item \textbf{Theorem}, \textbf{Proposition}, \textbf{Lemma},
\textbf{Corollary}: proved statements, in roughly decreasing order of
importance.
\item \textbf{Conjecture}: a statement that is not known to be true.
\item \textbf{Example}: an instance of a definition or a result.
\item \textbf{Counterexample}: shows that a hypothesis cannot be dropped, or
that a plausible statement is false.
\item \textbf{Remark}: context, motivation or intuition. Remarks are never used
in proofs.
\item \textbf{Notation}: fixes notation or a convention.
\item \textbf{Exercise}: a statement for the reader to prove, sometimes followed
by a solution.
\end{itemize}
\end{notation}

\begin{remark}\label{introduction:remark-proofs}
A result followed by \emph{Proof} has a complete proof that uses only earlier
material. A result followed by \emph{Sketch of proof} has an argument whose
main ideas are given but whose omitted verifications are named. A result
followed by no proof is quoted from the literature, with a reference.
Sketches and quotations are temporary: the aim of the project is to replace
all of them with complete proofs.
\end{remark}

\begin{remark}\label{introduction:remark-no-forward-references}
The project has no forward references: a proof may only use definitions and
results that appear before it. Examples, counterexamples, remarks and
exercises may mention notions that are defined later (for instance, the
category of topological spaces appears as an example in
Chapter~\ref{categories}), but such mentions are never used in a proof. The
only exception is this introduction, which refers forward freely.
\end{remark}

\begin{remark}\label{introduction:remark-tags}
Every labelled item (chapter, section, definition, result, example, and so on)
carries a permanent four-character \emph{tag}, shown next to it on the
website. The number of an item changes whenever material is inserted before
it; its tag never changes. To refer to an item from outside the project, cite
its tag.
\end{remark}

\begin{remark}\label{introduction:remark-status}
Each chapter has a status. \emph{Planned}: only an outline exists.
\emph{Draft}: the chapter is being written, and its proofs have not been
checked. \emph{Complete}: all planned material is written, and the proofs
await independent checking. \emph{Checked}: every proof in the chapter has been
checked by a reader other than its author.
\end{remark}

\section{Mathematical conventions}\label{introduction:section-mathematical-conventions}

\begin{notation}\label{introduction:notation-universe}
We work in ZFC set theory together with the assumption that a Grothendieck
universe exists, and we fix such a universe $\mathcal{U}$ once and for all. A
set is \emph{small} if it is an element of $\mathcal{U}$. Chapter~\ref{sets}
develops these notions; Chapter~\ref{categories} explains why they are needed.
\end{notation}

\begin{notation}\label{introduction:notation-numbers}
The natural numbers include zero: $\NN = \{0, 1, 2, \ldots\}$. This is the
convention of set theory and logic, where $\NN$ is built starting from
$0 = \emptyset$ and each natural number $n = \{0, 1, \ldots, n - 1\}$ is a set with
$n$ elements (Remark~\ref{sets:remark-zero-natural}). Many texts, especially in
analysis and number theory, start from $1$ instead; for the positive integers we
write $\NN_{>0} = \{1, 2, 3, \ldots\}$, and more generally
$\NN_{\geq k} = \{k, k + 1, \ldots\}$. A real number $x$ is \emph{positive} if
$x > 0$ and \emph{non-negative} if $x \geq 0$.
\end{notation}

\begin{notation}\label{introduction:notation-maps}
The composition of maps $f \colon X \to Y$ and $g \colon Y \to Z$ is written
$g \circ f$ or $gf$. The identity map of $X$ is written $\id_X$.
\end{notation}

\begin{notation}\label{introduction:notation-rings}
A ring is associative and has an identity element $1$, and ring homomorphisms
preserve identity elements. Rings are not assumed to be commutative, except
from Chapter~\ref{commutative-algebra} onwards, where ``ring'' means
``commutative ring'' unless stated otherwise.
\end{notation}
